\definecolor[darkgray] [s=.125]
\definecolor[lightgray][s=.875]

\setupcolors
  [state=start,
   textcolor=lightgray]

\newbox\foxetbox
\setbox\foxetbox\hbox{\chardef\foxetcolormode2\scale[width=400pt]{\strut\foXet}}

\definepapersize
  [foxet]
  [width=\wd\foxetbox,
   height=\dimexpr(\ht\foxetbox+\dp\foxetbox)]

\setuppapersize
  [foxet]
  [foxet]

\setupbackgrounds
  [page]
  [background=color,
   backgroundcolor=darkgray]

\setuplayout
  [backspace=8pt,
   topspace=8pt,
   header=0pt,
   footer=0pt,
   width=middle,
   height=middle]

\setupwhitespace
  [halfline]

\starttext

\startTEXpage\copy\foxetbox\stopTEXpage

\definedfont[cmvtt10]

You can process a file containing so called XML formatting objects
by issuing the command:

\starttyping
texexec --pdf --use=foxet yourfile.fo
texexec --pdf --use=foxet yourfile.xml
texexec --pdf --use=foxet yourfile.xyz --xml
\stoptyping

The shortcut \type {--foxet} can be used instead:

\starttyping
texexec --foxet --pdf yourfile.fox
\stoptyping

If your file does not need graphics and/or cross references, you
can do with one run and gain some speed by using the \type {--once}
directive.

\vfill

\hfill Hans Hagen, PRAGMA-ADE, 2004

\stoptext